%%=============================================================================
%% Voorwoord
%%=============================================================================

\chapter*{\IfLanguageName{dutch}{Woord vooraf}{Preface}}%
\label{ch:voorwoord}

%% TODO:
%% Het voorwoord is het enige deel van de bachelorproef waar je vanuit je
%% eigen standpunt (``ik-vorm'') mag schrijven. Je kan hier bv. motiveren
%% waarom jij het onderwerp wil bespreken.
%% Vergeet ook niet te bedanken wie je geholpen/gesteund/... heeft

%%=============================================================================
%% Voorwoord
%%=============================================================================


Tijdens mijn studie kwam ik voor het eerst in contact met mainframe-technologie. Hoewel deze wereld in het begin wat onbekend leek, raakte ik al snel gefascineerd door de enorme capaciteit en de performance van deze systemen. Het feit dat mainframes de ruggengraat vormen van de grootste organisaties ter wereld, gaf mij de motivatie om hier meer over te willen weten. 

Tijdens dit proces ontdekte ik HLASM, de assemblytaal van IBM. Mijn interesse werd meteen gewekt omdat deze taal je dwingt om na te denken over de directe werking van de hardware. Mijn overtuiging is dat het leren van een assemblytaal je een betere programmeur maakt. Door op het laagste niveau te werken, krijg je een veel dieper inzicht in hoe instructies en het geheugen echt functioneren. Dit is fundamentele kennis die vaak verloren gaat bij het gebruik van moderne programmeertalen.

HLASM heeft echter een steile leercurve, wat het voor beginners lastig maakt. Omdat ik zelf ook nog volop aan het leren ben, kon ik mij goed verplaatsen in de uitdagingen waar een nieuwe student tegenaan loopt. Vanuit die gedachte ontstond het idee voor deze bachelorproef: het bouwen van een tool die de control-flow van HLASM-programma's visueel voorstelt om het leerproces te ondersteunen.

Het uitwerken van dit project heeft mij niet alleen technische kennis opgeleverd, maar ook inzicht gegeven in hoe complexe informatie vereenvoudigd kan worden. Het bouwen van de proof-of-concept was een waardevolle kans om theorie om te zetten in iets concreets.

\paragraph{Dankwoord}

Graag wil ik mijn promotor, L. Blondeel, bedanken voor de begeleiding en de nuttige feedback tijdens het schrijven van deze bachelorproef.

Daarnaast bedank ik mijn co-promotor, B. Tjassens Keiser, voor de hulp bij technische vragen.

Tot slot wil ik mijn familie en vrienden bedanken voor hun steun en motivatie. Hun aanmoedigingen hebben mij geholpen om dit traject succesvol af te ronden.